\chapter{Unit 2: Problem Definition and Analysis}


\section{Section A: Multiple Choice Practice Questions (60 Curated Problems)}

\noindent Practice and test your mastery across all core concepts of this unit. Each question reflects standard university examination patterns and Bloom's Taxonomy competencies.
\vspace{3mm}

\noindent\textbf{Q1. Who is a stakeholder in a design problem?} \hfill [Tag: Understanding stakeholder needs and contextual factors influencing problems | Easy]
\begin{enumerate}[label=(\alph*),topsep=2pt,itemsep=1pt,partopsep=0pt,parsep=0pt]
    \item {A tool used to solve the problem}
    \item {A person affected by the problem}
    \item {A final version of the product}
    \item {A person unrelated to the problem}
\end{enumerate}
\vspace{2mm}

\noindent\textbf{Q2. What is a stakeholder need?} \hfill [Tag: Understanding stakeholder needs and contextual factors influencing problems | Easy]
\begin{enumerate}[label=(\alph*),topsep=2pt,itemsep=1pt,partopsep=0pt,parsep=0pt]
    \item {A feature chosen by the designer}
    \item {A requirement or expectation of an affected person}
    \item {A list of available design tools}
    \item {A limit placed on the project budget}
\end{enumerate}
\vspace{2mm}

\noindent\textbf{Q3. Which factor describes the setting in which a problem occurs?} \hfill [Tag: Understanding stakeholder needs and contextual factors influencing problems | Easy]
\begin{enumerate}[label=(\alph*),topsep=2pt,itemsep=1pt,partopsep=0pt,parsep=0pt]
    \item {Prototype}
    \item {Deliverable}
    \item {Context}
    \item {Objective}
\end{enumerate}
\vspace{2mm}

\noindent\textbf{Q4. Why should designers consider different stakeholders?} \hfill [Tag: Understanding stakeholder needs and contextual factors influencing problems | Easy]
\begin{enumerate}[label=(\alph*),topsep=2pt,itemsep=1pt,partopsep=0pt,parsep=0pt]
    \item {To guarantee one perfect solution}
    \item {To avoid collecting information}
    \item {To remove all project constraints}
    \item {To understand different needs}
\end{enumerate}
\vspace{2mm}

\noindent\textbf{Q5. Which technique involves asking people questions to learn about a problem?} \hfill [Tag: Techniques for problem identification | Easy]
\begin{enumerate}[label=(\alph*),topsep=2pt,itemsep=1pt,partopsep=0pt,parsep=0pt]
    \item {Prototyping}
    \item {Benchmarking}
    \item {Interviewing}
    \item {Scheduling}
\end{enumerate}
\vspace{2mm}

\noindent\textbf{Q6. What is observation used for during problem identification?} \hfill [Tag: Techniques for problem identification | Easy]
\begin{enumerate}[label=(\alph*),topsep=2pt,itemsep=1pt,partopsep=0pt,parsep=0pt]
    \item {Selecting a project budget}
    \item {Writing the final project report}
    \item {Watching actions in context}
    \item {Creating a product advertisement}
\end{enumerate}
\vspace{2mm}

\noindent\textbf{Q7. What is the main purpose of brainstorming during problem identification?} \hfill [Tag: Techniques for problem identification | Easy]
\begin{enumerate}[label=(\alph*),topsep=2pt,itemsep=1pt,partopsep=0pt,parsep=0pt]
    \item {Measuring product weight}
    \item {Approving the project budget}
    \item {Selecting one final answer}
    \item {Generating many ideas}
\end{enumerate}
\vspace{2mm}

\noindent\textbf{Q8. Which tool is commonly used to identify the causes of a problem?} \hfill [Tag: Techniques for problem identification | Easy]
\begin{enumerate}[label=(\alph*),topsep=2pt,itemsep=1pt,partopsep=0pt,parsep=0pt]
    \item {Customer invoice}
    \item {Sales receipt}
    \item {Project calendar}
    \item {Cause-and-effect diagram}
\end{enumerate}
\vspace{2mm}

\noindent\textbf{Q9. What does the repeated use of the question "Why?" help reveal?} \hfill [Tag: Techniques for problem identification | Easy]
\begin{enumerate}[label=(\alph*),topsep=2pt,itemsep=1pt,partopsep=0pt,parsep=0pt]
    \item {Team attendance}
    \item {Available office space}
    \item {Final product colors}
    \item {Underlying causes}
\end{enumerate}
\vspace{2mm}

\noindent\textbf{Q10. What is the purpose of a problem statement?} \hfill [Tag: Problem statement formulation | Easy]
\begin{enumerate}[label=(\alph*),topsep=2pt,itemsep=1pt,partopsep=0pt,parsep=0pt]
    \item {To list every possible design}
    \item {To report the team's expenses}
    \item {To clearly describe the problem}
    \item {To advertise the proposed solution}
\end{enumerate}
\vspace{2mm}

\noindent\textbf{Q11. Which statement is the best example of a problem statement?} \hfill [Tag: Problem statement formulation | Easy]
\begin{enumerate}[label=(\alph*),topsep=2pt,itemsep=1pt,partopsep=0pt,parsep=0pt]
    \item {We will build a mobile library application.}
    \item {The project team will meet every Friday.}
    \item {The application will use blue and white colors.}
    \item {Students need easier access to library resources.}
\end{enumerate}
\vspace{2mm}

\noindent\textbf{Q12. A well-written problem statement should usually be:} \hfill [Tag: Problem statement formulation | Easy]
\begin{enumerate}[label=(\alph*),topsep=2pt,itemsep=1pt,partopsep=0pt,parsep=0pt]
    \item {Limited to a proposed technology}
    \item {Long and unrelated}
    \item {Based only on assumptions}
    \item {Clear and focused}
\end{enumerate}
\vspace{2mm}

\noindent\textbf{Q13. Which item should generally be avoided in a problem statement?} \hfill [Tag: Problem statement formulation | Easy]
\begin{enumerate}[label=(\alph*),topsep=2pt,itemsep=1pt,partopsep=0pt,parsep=0pt]
    \item {The affected user group}
    \item {A predetermined solution}
    \item {Relevant background information}
    \item {The main difficulty}
\end{enumerate}
\vspace{2mm}

\noindent\textbf{Q14. What does project scope define?} \hfill [Tag: Scope definition | Easy]
\begin{enumerate}[label=(\alph*),topsep=2pt,itemsep=1pt,partopsep=0pt,parsep=0pt]
    \item {The order of employee birthdays}
    \item {The color of every document}
    \item {The team's personal preferences}
    \item {What the project includes}
\end{enumerate}
\vspace{2mm}

\noindent\textbf{Q15. What is an out-of-scope item?} \hfill [Tag: Scope definition | Easy]
\begin{enumerate}[label=(\alph*),topsep=2pt,itemsep=1pt,partopsep=0pt,parsep=0pt]
    \item {An item placed at the project center}
    \item {An item completed before planning}
    \item {An item required by every stakeholder}
    \item {An item excluded from the project}
\end{enumerate}
\vspace{2mm}

\noindent\textbf{Q16. Why is defining scope important?} \hfill [Tag: Scope definition | Easy]
\begin{enumerate}[label=(\alph*),topsep=2pt,itemsep=1pt,partopsep=0pt,parsep=0pt]
    \item {It guarantees that no changes occur}
    \item {It prevents uncontrolled expansion}
    \item {It makes every stakeholder agree immediately}
    \item {It eliminates the need for research}
\end{enumerate}
\vspace{2mm}

\noindent\textbf{Q17. What is a constraint?} \hfill [Tag: Constraint analysis | Easy]
\begin{enumerate}[label=(\alph*),topsep=2pt,itemsep=1pt,partopsep=0pt,parsep=0pt]
    \item {A creative design idea}
    \item {A method for interviewing users}
    \item {A limit on a project}
    \item {A possible customer benefit}
\end{enumerate}
\vspace{2mm}

\noindent\textbf{Q18. Which is an example of a time constraint?} \hfill [Tag: Constraint analysis | Easy]
\begin{enumerate}[label=(\alph*),topsep=2pt,itemsep=1pt,partopsep=0pt,parsep=0pt]
    \item {The product may help reduce confusion.}
    \item {The solution must be ready in four weeks.}
    \item {The team wants to study user habits.}
    \item {The solution should improve user comfort.}
\end{enumerate}
\vspace{2mm}

\noindent\textbf{Q19. Which of the following is a common project constraint?} \hfill [Tag: Constraint analysis | Easy]
\begin{enumerate}[label=(\alph*),topsep=2pt,itemsep=1pt,partopsep=0pt,parsep=0pt]
    \item {Motivation}
    \item {Imagination}
    \item {Curiosity}
    \item {Budget}
\end{enumerate}
\vspace{2mm}

\noindent\textbf{Q20. What is an objective?} \hfill [Tag: Objective setting | Easy]
\begin{enumerate}[label=(\alph*),topsep=2pt,itemsep=1pt,partopsep=0pt,parsep=0pt]
    \item {A stakeholder interview}
    \item {A description of a past failure}
    \item {A desired result}
    \item {A project limitation}
\end{enumerate}
\vspace{2mm}

\noindent\textbf{Q21. A hospital wants to reduce patient waiting times. Patients value speed, nurses value manageable workloads, and administrators value cost control. Which approach best addresses the different stakeholder needs?} \hfill [Tag: Understanding stakeholder needs and contextual factors influencing problems | Medium]
\begin{enumerate}[label=(\alph*),topsep=2pt,itemsep=1pt,partopsep=0pt,parsep=0pt]
    \item {Map each stakeholder's needs and identify important trade-offs}
    \item {Select the need that can be measured most easily}
    \item {Average all stakeholder preferences into one target}
    \item {Prioritize the needs of the hospital administrator}
\end{enumerate}
\vspace{2mm}

\noindent\textbf{Q22. A public transportation app has low usage in a rural area. Which contextual factor should be investigated first?} \hfill [Tag: Understanding stakeholder needs and contextual factors influencing problems | Medium]
\begin{enumerate}[label=(\alph*),topsep=2pt,itemsep=1pt,partopsep=0pt,parsep=0pt]
    \item {The reliability of mobile internet coverage}
    \item {The number of app features offered}
    \item {The average age of the design team}
    \item {The color scheme of the app interface}
\end{enumerate}
\vspace{2mm}

\noindent\textbf{Q23. During interviews, teachers request more digital tools, but students report difficulty using existing tools. What is the most appropriate interpretation?} \hfill [Tag: Understanding stakeholder needs and contextual factors influencing problems | Medium]
\begin{enumerate}[label=(\alph*),topsep=2pt,itemsep=1pt,partopsep=0pt,parsep=0pt]
    \item {Teachers are the only stakeholders who matter}
    \item {Students are resisting all forms of technology}
    \item {More tools should be purchased immediately}
    \item {The request may reflect a need for simpler and better-supported tools}
\end{enumerate}
\vspace{2mm}

\noindent\textbf{Q24. A team notices that customers abandon an online payment process. Which technique would best reveal where users encounter difficulty?} \hfill [Tag: Techniques for problem identification | Medium]
\begin{enumerate}[label=(\alph*),topsep=2pt,itemsep=1pt,partopsep=0pt,parsep=0pt]
    \item {Creating a stakeholder power grid}
    \item {Mapping the customer journey}
    \item {Writing a broad vision statement}
    \item {Listing the team's available skills}
\end{enumerate}
\vspace{2mm}

\noindent\textbf{Q25. A team repeatedly asks "Why?" after identifying that deliveries are late. What is the main purpose of this technique?} \hfill [Tag: Techniques for problem identification | Medium]
\begin{enumerate}[label=(\alph*),topsep=2pt,itemsep=1pt,partopsep=0pt,parsep=0pt]
    \item {To uncover possible underlying causes}
    \item {To assign responsibility for the delay}
    \item {To rank solutions by implementation cost}
    \item {To determine which stakeholder has authority}
\end{enumerate}
\vspace{2mm}

\noindent\textbf{Q26. A school assumes that students need longer library hours because attendance is low. Which action would best test this assumption?} \hfill [Tag: Techniques for problem identification | Medium]
\begin{enumerate}[label=(\alph*),topsep=2pt,itemsep=1pt,partopsep=0pt,parsep=0pt]
    \item {Extend the hours for the entire semester}
    \item {Compare the library with an unrelated facility}
    \item {Ask the principal to select a new schedule}
    \item {Interview students and observe library use patterns}
\end{enumerate}
\vspace{2mm}

\noindent\textbf{Q27. A cause-and-effect diagram shows that a manufacturing defect may result from materials, machines, methods, or people. How should the team use this diagram?} \hfill [Tag: Techniques for problem identification | Medium]
\begin{enumerate}[label=(\alph*),topsep=2pt,itemsep=1pt,partopsep=0pt,parsep=0pt]
    \item {Replace data collection with team judgment}
    \item {Use it to organize hypotheses for further investigation}
    \item {Select the cause that appears at the top}
    \item {Treat every listed cause as equally proven}
\end{enumerate}
\vspace{2mm}

\noindent\textbf{Q28. Which problem statement is most suitable for a design thinking project?} \hfill [Tag: Problem statement formulation | Medium]
\begin{enumerate}[label=(\alph*),topsep=2pt,itemsep=1pt,partopsep=0pt,parsep=0pt]
    \item {Hospital staff need a faster computer system}
    \item {How might we help patients schedule appointments with less uncertainty?}
    \item {The hospital must purchase a better scheduling platform}
    \item {Create a mobile app for hospital appointments}
\end{enumerate}
\vspace{2mm}

\noindent\textbf{Q29. A team writes, "Employees are lazy and do not complete training." Which revision is more useful and less biased?} \hfill [Tag: Problem statement formulation | Medium]
\begin{enumerate}[label=(\alph*),topsep=2pt,itemsep=1pt,partopsep=0pt,parsep=0pt]
    \item {Employees should be replaced if training is incomplete}
    \item {Employees should receive more reminders from management}
    \item {Employees need clearer access and time to complete required training}
    \item {Employees need stricter penalties for missed training}
\end{enumerate}
\vspace{2mm}

\noindent\textbf{Q30. Which element is most important when converting a general issue into a focused problem statement?} \hfill [Tag: Problem statement formulation | Medium]
\begin{enumerate}[label=(\alph*),topsep=2pt,itemsep=1pt,partopsep=0pt,parsep=0pt]
    \item {A clearly identified user and unmet need}
    \item {A detailed list of team responsibilities}
    \item {A complete implementation budget}
    \item {A preferred product or technology}
\end{enumerate}
\vspace{2mm}

\noindent\textbf{Q31. A team changes "How can we reduce complaints?" to "How might we help first-time customers understand return policies before purchase?" What improvement does the revision make?} \hfill [Tag: Problem statement formulation | Medium]
\begin{enumerate}[label=(\alph*),topsep=2pt,itemsep=1pt,partopsep=0pt,parsep=0pt]
    \item {It identifies a user, context, and specific opportunity}
    \item {It focuses only on reducing the organization's workload}
    \item {It guarantees that an information campaign will work}
    \item {It removes the need to study customer behavior}
\end{enumerate}
\vspace{2mm}

\noindent\textbf{Q32. A project aims to improve food delivery in a university. Which scope is the most manageable for an initial study?} \hfill [Tag: Scope definition | Medium]
\begin{enumerate}[label=(\alph*),topsep=2pt,itemsep=1pt,partopsep=0pt,parsep=0pt]
    \item {Improve all food systems across the entire city}
    \item {Solve every problem faced by university students}
    \item {Redesign food delivery for one campus dining hall}
    \item {Replace all local food delivery companies}
\end{enumerate}
\vspace{2mm}

\noindent\textbf{Q33. A project team includes improving parking, public transport, classroom design, and cafeteria service, although its original goal concerns classroom navigation. What scope problem is present?} \hfill [Tag: Scope definition | Medium]
\begin{enumerate}[label=(\alph*),topsep=2pt,itemsep=1pt,partopsep=0pt,parsep=0pt]
    \item {The project has excessive scope expansion}
    \item {The project has already validated its solution}
    \item {The project has no possible constraints}
    \item {The project has too few stakeholders}
\end{enumerate}
\vspace{2mm}

\noindent\textbf{Q34. Which statement best defines an out-of-scope item?} \hfill [Tag: Scope definition | Medium]
\begin{enumerate}[label=(\alph*),topsep=2pt,itemsep=1pt,partopsep=0pt,parsep=0pt]
    \item {A related issue intentionally excluded from the current project}
    \item {A user need that should guide research}
    \item {A measurable result used to judge success}
    \item {A constraint that limits the available budget}
\end{enumerate}
\vspace{2mm}

\noindent\textbf{Q35. A nonprofit has only two months and a limited budget to improve access to legal information. How should these conditions be treated?} \hfill [Tag: Constraint analysis | Medium]
\begin{enumerate}[label=(\alph*),topsep=2pt,itemsep=1pt,partopsep=0pt,parsep=0pt]
    \item {As constraints that shape feasible project choices}
    \item {As stakeholder needs that define the user group}
    \item {As evidence that the problem should be ignored}
    \item {As objectives that must be maximized}
\end{enumerate}
\vspace{2mm}

\noindent\textbf{Q36. A proposed healthcare solution requires collecting sensitive patient data. Which constraint should receive particular attention?} \hfill [Tag: Constraint analysis | Medium]
\begin{enumerate}[label=(\alph*),topsep=2pt,itemsep=1pt,partopsep=0pt,parsep=0pt]
    \item {Privacy and regulatory requirements}
    \item {The color of the project timeline}
    \item {The team's preferred presentation style}
    \item {The number of brainstorming sessions}
\end{enumerate}
\vspace{2mm}

\noindent\textbf{Q37. A team lists low funding, limited staff, safety regulations, and unreliable internet access. What is the best next step?} \hfill [Tag: Constraint analysis | Medium]
\begin{enumerate}[label=(\alph*),topsep=2pt,itemsep=1pt,partopsep=0pt,parsep=0pt]
    \item {Ignore constraints until the final design}
    \item {Assume every constraint has the same importance}
    \item {Remove constraints that make the project difficult}
    \item {Prioritize constraints by impact and negotiability}
\end{enumerate}
\vspace{2mm}

\noindent\textbf{Q38. Which objective is most measurable for a project improving appointment scheduling?} \hfill [Tag: Objective setting | Medium]
\begin{enumerate}[label=(\alph*),topsep=2pt,itemsep=1pt,partopsep=0pt,parsep=0pt]
    \item {Reduce average scheduling time by 20\% within three months}
    \item {Ensure that users feel completely satisfied}
    \item {Develop the most innovative booking experience}
    \item {Make scheduling much better for everyone}
\end{enumerate}
\vspace{2mm}

\noindent\textbf{Q39. A team wants to improve access to educational materials. Which objective best balances outcome and feasibility?} \hfill [Tag: Objective setting | Medium]
\begin{enumerate}[label=(\alph*),topsep=2pt,itemsep=1pt,partopsep=0pt,parsep=0pt]
    \item {Make the resource system highly successful}
    \item {Eliminate all educational inequality within one year}
    \item {Increase weekly access to required materials among target learners by 15\%}
    \item {Provide every learner with every possible resource}
\end{enumerate}
\vspace{2mm}

\noindent\textbf{Q40. A project objective is to reduce errors in a warehouse process. Which indicator would best measure achievement?} \hfill [Tag: Objective setting | Medium]
\begin{enumerate}[label=(\alph*),topsep=2pt,itemsep=1pt,partopsep=0pt,parsep=0pt]
    \item {The length of the final project report}
    \item {The number of brainstorming ideas generated}
    \item {The percentage of orders processed without errors}
    \item {The number of project meetings completed}
\end{enumerate}
\vspace{2mm}

\noindent\textbf{Q41. A hospital introduces a digital check-in system. Patients report shorter queues, reception staff report increased workload, and clinicians report missing allergy information. Which interpretation best reflects the stakeholder analysis?} \hfill [Tag: Understanding stakeholder needs and contextual factors influencing problems | Hard]
\begin{enumerate}[label=(\alph*),topsep=2pt,itemsep=1pt,partopsep=0pt,parsep=0pt]
    \item {The system should prioritize clinician preferences exclusively}
    \item {The system fails because reception staff workload increased}
    \item {The system has uneven effects across stakeholder outcomes}
    \item {The system succeeds because average waiting time decreased}
\end{enumerate}
\vspace{2mm}

\noindent\textbf{Q42. A public-transport redesign is based mainly on interviews with frequent commuters. Later, occasional users and people with disabilities report severe navigation difficulties. Which analytical error most directly caused this result?} \hfill [Tag: Understanding stakeholder needs and contextual factors influencing problems | Hard]
\begin{enumerate}[label=(\alph*),topsep=2pt,itemsep=1pt,partopsep=0pt,parsep=0pt]
    \item {Treating stated preferences as measurable outcomes}
    \item {Confusing operational constraints with user needs}
    \item {Sampling visible users rather than affected users}
    \item {Using qualitative evidence before quantitative evidence}
\end{enumerate}
\vspace{2mm}

\noindent\textbf{Q43. A school wants to reduce absenteeism. Parents cite transportation, teachers cite disengagement, and students cite fear of bullying. Which next step is most defensible before selecting an intervention?} \hfill [Tag: Understanding stakeholder needs and contextual factors influencing problems | Hard]
\begin{enumerate}[label=(\alph*),topsep=2pt,itemsep=1pt,partopsep=0pt,parsep=0pt]
    \item {Map causal relationships and test assumptions with each group}
    \item {Combine all explanations into one broad attendance objective}
    \item {Rank the explanations by the number of people mentioning them}
    \item {Select the explanation that administrators can address fastest}
\end{enumerate}
\vspace{2mm}

\noindent\textbf{Q44. A company observes that customer complaints about billing have doubled, although the number of customers has also doubled. Which technique best prevents premature labeling of this observation as a worsening problem?} \hfill [Tag: Techniques for problem identification | Hard]
\begin{enumerate}[label=(\alph*),topsep=2pt,itemsep=1pt,partopsep=0pt,parsep=0pt]
    \item {Treat the absolute complaint count as the primary indicator}
    \item {Compare complaint rates per customer across equivalent periods}
    \item {Interview only customers who submitted recent complaints}
    \item {Define the problem as insufficient customer-service staffing}
\end{enumerate}
\vspace{2mm}

\noindent\textbf{Q45. During a root-cause investigation, the team repeatedly asks why a medical sample was mislabeled and eventually identifies an unclear handoff protocol. What is the main limitation of relying only on this technique?} \hfill [Tag: Techniques for problem identification | Hard]
\begin{enumerate}[label=(\alph*),topsep=2pt,itemsep=1pt,partopsep=0pt,parsep=0pt]
    \item {It cannot reveal any process-related causes}
    \item {It requires numerical data before questions can be asked}
    \item {It may impose a single linear cause on a systemic problem}
    \item {It prevents participants from proposing possible solutions}
\end{enumerate}
\vspace{2mm}

\noindent\textbf{Q46. A service has acceptable average performance, but a small subgroup experiences extreme delays. Which investigation approach is most appropriate?} \hfill [Tag: Techniques for problem identification | Hard]
\begin{enumerate}[label=(\alph*),topsep=2pt,itemsep=1pt,partopsep=0pt,parsep=0pt]
    \item {Remove outliers before calculating service performance}
    \item {Segment the data to identify distributional inequities}
    \item {Replace quantitative analysis with a general satisfaction survey}
    \item {Optimize the mean because it represents typical performance}
\end{enumerate}
\vspace{2mm}

\noindent\textbf{Q47. Which problem statement is most suitable for design exploration?} \hfill [Tag: Problem statement formulation | Hard]
\begin{enumerate}[label=(\alph*),topsep=2pt,itemsep=1pt,partopsep=0pt,parsep=0pt]
    \item {The appointment system must be replaced before the next quarter}
    \item {The clinic needs an automated appointment application}
    \item {Patients fail to attend because staff lack discipline}
    \item {How might the clinic help patients complete appointments reliably?}
\end{enumerate}
\vspace{2mm}

\noindent\textbf{Q48. A team defines its problem as: "Small retailers need better inventory software." Which revision most improves the statement while preserving analytical neutrality?} \hfill [Tag: Problem statement formulation | Hard]
\begin{enumerate}[label=(\alph*),topsep=2pt,itemsep=1pt,partopsep=0pt,parsep=0pt]
    \item {Small retailers struggle to maintain stock availability under uncertain demand}
    \item {Small retailers should purchase software from the market leader}
    \item {Small retailers need artificial-intelligence inventory software}
    \item {Small retailers require a dashboard with real-time visualizations}
\end{enumerate}
\vspace{2mm}

\noindent\textbf{Q49. A project team writes, "Employees are unmotivated, causing low productivity." Which issue is most important to resolve before adopting this as the formal problem statement?} \hfill [Tag: Problem statement formulation | Hard]
\begin{enumerate}[label=(\alph*),topsep=2pt,itemsep=1pt,partopsep=0pt,parsep=0pt]
    \item {The statement is too narrow to permit any intervention}
    \item {The presumed cause is treated as established without evidence}
    \item {The statement focuses on employees rather than equipment}
    \item {The statement contains too many measurable variables}
\end{enumerate}
\vspace{2mm}

\noindent\textbf{Q50. A nonprofit wants to increase food access but risks defining the problem as "low food donations." Which formulation best avoids confusing an indicator with the underlying need?} \hfill [Tag: Problem statement formulation | Hard]
\begin{enumerate}[label=(\alph*),topsep=2pt,itemsep=1pt,partopsep=0pt,parsep=0pt]
    \item {How might we improve reliable access to nutritious food for households?}
    \item {How might we increase the number of donated food packages?}
    \item {How might we persuade more supermarkets to donate surplus food?}
    \item {How might we reduce the nonprofit's procurement expenses?}
\end{enumerate}
\vspace{2mm}

\noindent\textbf{Q51. A university investigates delayed graduation. Which scope boundary is most useful for an initial study intended to produce actionable findings within one semester?} \hfill [Tag: Scope definition | Hard]
\begin{enumerate}[label=(\alph*),topsep=2pt,itemsep=1pt,partopsep=0pt,parsep=0pt]
    \item {Only students who graduated on time during the previous decade}
    \item {All educational experiences across every university system}
    \item {Every social, financial, and academic factor affecting graduation}
    \item {Course scheduling in two high-enrollment departments this year}
\end{enumerate}
\vspace{2mm}

\noindent\textbf{Q52. A team studying emergency-department congestion excludes ambulance arrivals because they are difficult to influence. What is the principal risk of this scope decision?} \hfill [Tag: Scope definition | Hard]
\begin{enumerate}[label=(\alph*),topsep=2pt,itemsep=1pt,partopsep=0pt,parsep=0pt]
    \item {The findings will automatically generalize to other hospitals}
    \item {The project will necessarily exceed its budget}
    \item {The analysis may omit a major driver of observed congestion}
    \item {The team will be unable to collect any valid data}
\end{enumerate}
\vspace{2mm}

\noindent\textbf{Q53. A project initially targets reducing workplace injuries in one factory. Midway through, the sponsor requests inclusion of all regional factories and long-term employee wellness. What is the best response?} \hfill [Tag: Scope definition | Hard]
\begin{enumerate}[label=(\alph*),topsep=2pt,itemsep=1pt,partopsep=0pt,parsep=0pt]
    \item {Accept the additions without changing schedule or resources}
    \item {Assess impacts and renegotiate boundaries, resources, and outcomes}
    \item {Include the additions informally without documenting the decision}
    \item {Reject all additions because scope can never change}
\end{enumerate}
\vspace{2mm}

\noindent\textbf{Q54. A team can either study every stage of a customer journey superficially or study one high-loss stage deeply. Which principle should guide the choice?} \hfill [Tag: Scope definition | Hard]
\begin{enumerate}[label=(\alph*),topsep=2pt,itemsep=1pt,partopsep=0pt,parsep=0pt]
    \item {Choose the broadest scope because completeness is always superior}
    \item {Choose the scope with the fewest stakeholders to consult}
    \item {Choose the scope that balances causal relevance and decision usefulness}
    \item {Choose the narrowest scope regardless of causal relevance}
\end{enumerate}
\vspace{2mm}

\noindent\textbf{Q55. A redesign must comply with privacy law, use existing hardware, launch in six months, and serve users with limited connectivity. Which classification is most accurate?} \hfill [Tag: Constraint analysis | Hard]
\begin{enumerate}[label=(\alph*),topsep=2pt,itemsep=1pt,partopsep=0pt,parsep=0pt]
    \item {Privacy, hardware, time, and connectivity shape the feasible solution space}
    \item {Privacy and hardware are constraints; time and connectivity are objectives}
    \item {Only the launch date is a true constraint because it is measurable}
    \item {All four are preferences that can be traded away}
\end{enumerate}
\vspace{2mm}

\noindent\textbf{Q56. A team claims that a fixed budget is an absolute constraint, but stakeholders would accept a delayed launch for a higher budget. What does this reveal?} \hfill [Tag: Constraint analysis | Hard]
\begin{enumerate}[label=(\alph*),topsep=2pt,itemsep=1pt,partopsep=0pt,parsep=0pt]
    \item {The launch date is not part of the problem context}
    \item {The budget must be removed from all analysis}
    \item {The budget is irrelevant to the project definition}
    \item {The budget is a negotiable constraint with trade-offs}
\end{enumerate}
\vspace{2mm}

\noindent\textbf{Q57. A proposed solution meets all stated requirements but cannot be operated by staff during peak demand. Which overlooked constraint is most likely involved?} \hfill [Tag: Constraint analysis | Hard]
\begin{enumerate}[label=(\alph*),topsep=2pt,itemsep=1pt,partopsep=0pt,parsep=0pt]
    \item {The solution's color palette}
    \item {The number of alternative concepts generated}
    \item {The team's preference for quantitative research}
    \item {The system's operational capacity and usability}
\end{enumerate}
\vspace{2mm}

\noindent\textbf{Q58. A project objective states: "Make the service better for everyone." Which revision is strongest?} \hfill [Tag: Objective setting | Hard]
\begin{enumerate}[label=(\alph*),topsep=2pt,itemsep=1pt,partopsep=0pt,parsep=0pt]
    \item {Ensure users and staff feel completely satisfied}
    \item {Reduce median completion time by 20\% without lowering accessibility scores}
    \item {Install a faster platform before the next planning meeting}
    \item {Improve the service as much as possible}
\end{enumerate}
\vspace{2mm}

\noindent\textbf{Q59. A city sets an objective to minimize traffic congestion. A proposed redesign reduces car travel time but increases pedestrian injuries. Which objective structure would best expose this failure?} \hfill [Tag: Objective setting | Hard]
\begin{enumerate}[label=(\alph*),topsep=2pt,itemsep=1pt,partopsep=0pt,parsep=0pt]
    \item {Measure congestion only during the least busy period}
    \item {Replace all outcome measures with stakeholder opinions}
    \item {Use travel time as the only success metric}
    \item {Set a weighted objective with a minimum safety threshold}
\end{enumerate}
\vspace{2mm}

\noindent\textbf{Q60. A team must choose between objectives that reduce cost, increase equity, and improve reliability. No option dominates on all three dimensions. What is the most rigorous next step?} \hfill [Tag: Objective setting | Hard]
\begin{enumerate}[label=(\alph*),topsep=2pt,itemsep=1pt,partopsep=0pt,parsep=0pt]
    \item {Choose the easiest objective to measure quantitatively}
    \item {Elicit stakeholder priorities and analyze explicit trade-offs}
    \item {Select the objective favored by the project manager}
    \item {Combine the objectives without defining their relative importance}
\end{enumerate}
\vspace{2mm}


\newpage
\section{Section B: Comprehensive Subjective Problem with Detailed Analytical Rubric}

\begin{subjectiveq}{2}{1}
\textbf{Case Background:} An emergency outpatient department experiences severe overcrowding and patient dissatisfaction. Average waiting times have escalated to 210 minutes. Clinical staff complain of burnout, while hospital management threatens budget cuts due to declining patient throughput.

\textbf{Required Design Deliverables:}
\begin{enumerate}[label=(\alph*)]
    \item Perform a structured Root Cause Analysis utilizing both the 5-Whys methodology and an Ishikawa (Fishbone) diagram categorization (People, Process, Policy, Plant/Equipment).
    \item Construct a comprehensive Stakeholder Matrix classifying at least four distinct stakeholder groups by Power and Interest.
    \item Define the project scope by establishing an unambiguous In-Scope vs. Out-of-Scope boundary ledger and formulate two SMART project objectives.
\end{enumerate}
\end{subjectiveq}

\begin{solbox}
\textbf{1. Structured Root Cause Analysis (5-Whys \& Ishikawa Analysis):}\\
\textit{5-Whys Chain:}
\begin{itemize}
    \item \textbf{Why 1:} Why are ER waiting times exceeding 210 minutes? Because non-critical Category 4/5 patients occupy triage bays.
    \item \textbf{Why 2:} Why are non-critical patients presenting at the emergency ward? Because local outpatient clinics close at 4:00 PM.
    \item \textbf{Why 3:} Why do patients not wait for the next morning? Because they lack triage clarity and fear sudden symptom escalation.
    \item \textbf{Why 4:} Why is triage unable to divert them immediately? Because the hospital intake policy mandates full physician consultation before discharge.
    \item \textbf{Why 5 (Root Cause):} The system lacks an upfront algorithmic digital fast-track triage and extended community clinic tele-guidance network.
\end{itemize}

\textit{Ishikawa Dimension Breakdown:}
\begin{itemize}
    \item \textbf{People:} Clinical fatigue; triage nurse cognitive overload; patient anxiety.
    \item \textbf{Process:} Linear registration bottlenecks; batch paper chart routing; delayed diagnostic lab turnaround.
    \item \textbf{Policy:} Rigid intake liability protocols; absence of nurse-led discharge pathways for minor ailments.
    \item \textbf{Plant/Equipment:} Limited dedicated observation beds; poorly segregated waiting zones for infectious vs non-infectious cases.
\end{itemize}

\textbf{2. Stakeholder Power-Interest Classification:}\\
\begin{tabularx}{\linewidth}{|l|l|l|X|} \hline
\textbf{Stakeholder} & \textbf{Power} & \textbf{Interest} & \textbf{Engagement Strategy} \\ \hline
Emergency Attending Physicians & High & High & Key Partner: Co-design diagnostic workflow and fast-track protocols. \\ \hline
Hospital Executive Board & High & Medium & Keep Satisfied: Regular ROI, liability compliance, and throughput updates. \\ \hline
Outpatients \& Family Escorts & Low & High & Keep Informed: Provide clear visual queue wait-time dashboards and text alerts. \\ \hline
Third-Party Lab Vendors & Low & Low & Monitor: Enforce contractual SLA turnarounds on blood/radiology panels. \\ \hline
\end{tabularx}

\textbf{3. Scope Boundaries \& SMART Objectives:}\\
\textit{Scope Boundary Ledger:}
\begin{itemize}
    \item \textbf{In-Scope:} Arrival-to-triage workflow redesign; rapid-track kiosk prototyping; digital patient status board; nurse-led minor injury protocol.
    \item \textbf{Out-of-Scope:} Capital physical building expansion; surgical theatre operational restructuring; renegotiating nurse union contracts.
\end{itemize}

\textit{SMART Objectives:}
\begin{enumerate}[label=Objective \arabic*:]
    \item Reduce median non-critical (ESI Level 4 and 5) patient door-to-provider consultation time from 210 minutes to under 45 minutes within 90 days of intervention deployment.
    \item Elevate patient perceived communication and satisfaction scores from 52\% to $\ge 85\%$ as measured by post-visit automated surveys over a 6-month evaluation window.
\end{enumerate}
\end{solbox}

\begin{examinertip}
\textbf{Criteria for Full 10/10 Marks:}\\
1. \textbf{Rigorous Theoretical Framing (3 Marks):} Explicitly identify and define the formal model or framework applied (e.g., d.school stages, 5-Whys, MCDA, SCAMPER, Blueprinting, Minto SCQA). Avoid vague hand-waving.\\
2. \textbf{Contextual Application (4 Marks):} Directly connect every step to the specific facts, users, and constraints given in the case scenario.\\
3. \textbf{Actionable Deliverables \& Invariants (3 Marks):} Provide precise, formatted tables, matrices, or SMART criteria with clear quantitative bounds.
\end{examinertip}
